\chapter{Queue}
\label{chap:queue}

\section{Properties of queue}
\label{sec:properties of queue}

\begin{figure}[H]
  \centering
  \includegraphics[width=0.8\textwidth]{Figures/pimg4.png}
  \caption{Queue Visualization}
  \label{fig:pimg4-queue-visualization}
\end{figure}

\begin{enumerate}
  \item \textbf{First In First Out}
  \item Both end of a queue is open.
\end{enumerate}

\section{Application of Queue}
\label{sec:application-of-queue}
\begin{itemize}
  \item Hospital, Canteen queue
  \item \textbf{CPU scheduling} for processes
  \item Railway System for simulating train.
  \item Simulation
\end{itemize}

\section{Operations on Queue}
\label{sec:operations-on-queue}
\begin{itemize}
  \item \textbf{deque}: Take out element which came first.
  \item \textbf{enqueue}: Insert element at the last of the queue.
  \item \textbf{isEmpty}: Check if the queue is empty.
  \item \textbf{isFull}: is the queue full in terms of memory.
  \item \textbf{peek}: See what element will be dequed.
\end{itemize}

\section{Queue Implementation}
\label{sec:queue-implementation}
The following is the implementation of queue data structure in \textbf{Python}.
\begin{listing}[H]
  \inputminted{python}{Code/queue.py}
  \caption{Queue Implementation}
  \label{lst:queue-implementation}
\end{listing}

\section{Circular Queue}
\label{sec:circular-queue}

A circular queue is an extended version of a standard queue where the last
element is connected to the first, forming a circular structure. This design
efficiently utilizes the queue's storage space by allowing elements to wrap
around from the end to the beginning of the underlying array.

\begin{figure}[H]
  \centering
  \includegraphics[width=0.8\textwidth]{Figures/pimg5.png}
  \caption{Circular Queue}
  \label{fig:pimg5-circular-queue}
\end{figure}

\subsection{Key Characteristics and Components}
A circular queue typically uses:
\begin{itemize}
  \item An underlying array (or list in Python) of a fixed size (\textbf{maxSize} or \textbf{k} or \textbf{capacity}) to store the queue elements.
  \item Two pointers:
    \begin{itemize}
      \item \textbf{FRONT} (or \textbf{head}): Tracks the first element of the queue.
      \item \textbf{REAR} (or \textbf{tail}): Tracks the last element of the queue.
    \end{itemize}
  \item \textbf{Initialization}: Initially, both \textbf{FRONT} and \textbf{REAR} are often set to -1 to indicate an empty queue.
\end{itemize}

\subsection*{Operations}
The primary operations of a circular queue include:

\subsubsection*{Enqueue (Adding an element)}
\subsubsection*{Dequeue (Removing an element)}
\subsubsection*{Front (Peek)}
\subsubsection*{isFull}
\subsubsection*{isEmpty}

\section{Priority Queue}
\label{sec:priority-queue}
Each element has different priority. The item with highest priority is dequeued
first. Other than that, everything is same as queue.
